\section{总结与延伸}

本讲从 raw objects 出发，建立了完整的数据加工视角：Transformation 决定保留什么结构；Filtering 把质量写成 target 和 score；Deduplication 用 hashing 与集合近似控制重复；Mixing 把 token budget 分配给不同能力；Synthetic data 则把任务、环境和 verifier 变成新的数据来源。

五个可带走的判断是：

\begin{enumerate}
\item 数据管线每一步都会改变分布，必须做端到端评估。
\item Filter 的 target data 隐含价值观和 domain bias，阈值不是中性参数。
\item Dedup 的难点是定义“重复”，不是计算 hash。
\item Mixing 是训练目标的一部分，小规模最优比例未必能外推。
\item Synthetic data 的核心竞争力是 verifier 和 environment，而不是生成量。
\end{enumerate}

下一讲进入 mid/post-training：当数据已经变成 instruction、preference 和 safety behavior，SFT、PPO、DPO 怎样把它们转化为可控模型行为。

\subsection{拓展阅读}

建议结合官方 Lecture 14 source、DCLM、Dolma、RegMix、OpenThoughts、SWE-smith 与 SWE-rebench 的原始论文或项目页面阅读，重点比较它们如何定义 target、验证质量和记录 provenance。

\begin{knowledgebox}{复现实验：做一条 1B-token mini data pipeline}
选择两个 domain（例如 web 与 code）各准备一批 raw documents。固定 parser 后，依次记录过滤前后 token 数、exact/near dedup 删除率、每个 source 的 unique tokens、候选 mixture 权重与预计 epochs。至少保留一个未参与调参的 held-out slice，比较两套 threshold 或 mixture 在 per-domain loss、污染命中率和人工错误样本上的差异。最终交付不仅是一个训练文件，还应包含 manifest、parser/filter/dedup 版本、随机种子、失败样本和一页 data card；这样下一次模型行为变化才能沿 provenance 复盘。
\end{knowledgebox}

\end{document}
